\documentclass[12pt]{article}
\usepackage[utf8]{inputenc}
\usepackage[T1]{fontenc}
\usepackage[margin=1in]{geometry}
\usepackage{setspace}
\usepackage{xcolor}
\usepackage[colorlinks=true, urlcolor=blue!60!black, linkcolor=blue!60!black]{hyperref}
\setlength{\parindent}{2em}
\setlength{\parskip}{1.25em}
\renewcommand{\baselinestretch}{1.0}
\begin{document}
\begin{center}
{\LARGE \textit{Palko v. Connecticut}, 302 U.S. 319 (1937)\par}
\vspace{0.8em}
{\large Case Brief\par}
\end{center}
\vspace{1em}

\section{Memory Jogger}

Rejecting total incorporation of the Bill of Rights, Justice Cardozo articulated the ``ordered liberty'' standard for selective incorporation and held 8--1 that the Fourteenth Amendment's Due Process Clause did not bar Connecticut from appealing legal errors to retry a murder defendant and seek the death penalty---a specific holding later overruled in \textit{Benton v. Maryland} (1969) (\href{https://www.courtlistener.com/opinion/102879/palko-v-connecticut/}{CourtListener}; \href{https://www.oyez.org/cases/1900-1940/302us319}{Oyez}).

\section{Detailed Case Facts}

On September 30, 1935, Frank Palka (whose surname was misspelled ``Palko'' in court records) broke into Gilman's Music Store on Fairfield Avenue in Bridgeport, Connecticut, stole a radio, and fatally shot two Bridgeport police officers, including Sergeant Thomas Kearney, while fleeing (\href{https://ballotpedia.org/Palko_v._Connecticut}{Ballotpedia}; \href{https://law.justia.com/cases/connecticut/supreme-court/1936/121-conn-669.html}{Justia Connecticut}). A Fairfield County grand jury indicted Palko for first-degree murder (\href{https://supreme.justia.com/cases/federal/us/302/319/}{Justia}).

At his first trial in Connecticut Superior Court, the trial judge excluded testimony regarding Palko's confession, restricted the prosecution's cross-examination of Palko to impeach his credibility, and instructed the jury on the distinction between first- and second-degree murder (\href{https://supreme.justia.com/cases/federal/us/302/319/}{Justia}). The jury convicted Palko of second-degree murder, and the court sentenced him to life imprisonment (\href{https://supreme.justia.com/cases/federal/us/302/319/}{Justia}). Relying on an 1886 Connecticut statute, Conn. Gen. Stat. \S~6494---which authorized the State, with the presiding trial judge's permission, to appeal rulings on questions of law in criminal cases ``in the same manner and to the same effect as if made by the accused''---prosecutors appealed to the Connecticut Supreme Court of Errors (\href{https://supreme.justia.com/cases/federal/us/302/319/}{Justia}).

The Connecticut Supreme Court of Errors reversed the second-degree murder judgment and ordered a new trial, holding that the trial court committed legal errors prejudicial to the State in excluding the confession, barring impeachment cross-examination, and misinstructing the jury (\href{https://supreme.justia.com/cases/federal/us/302/319/}{Justia}; \href{https://law.justia.com/cases/connecticut/supreme-court/1936/121-conn-669.html}{Justia Connecticut}). Before a jury was impaneled for the retrial, Palko objected that a second trial placed him twice in jeopardy for the same offense in violation of the Fourteenth Amendment (\href{https://supreme.justia.com/cases/federal/us/302/319/}{Justia}). The trial court overruled the objection, the second jury convicted Palko of first-degree murder, and the court sentenced him to death (\href{https://supreme.justia.com/cases/federal/us/302/319/}{Justia}).

\section{Procedural History}

Following the first trial in Fairfield County Superior Court (second-degree murder conviction and life sentence), the State obtained the trial judge's permission under Conn. Gen. Stat. \S~6494 and appealed to the Connecticut Supreme Court of Errors (\href{https://supreme.justia.com/cases/federal/us/302/319/}{Justia}). On July 30, 1936, in \textit{State v. Palko}, 121 Conn. 669, 186 A. 657 (1936), the Supreme Court of Errors reversed and ordered a new trial (\href{https://law.justia.com/cases/connecticut/supreme-court/1936/121-conn-669.html}{Justia Connecticut}).

On retrial in Superior Court, Palko's Fourteenth Amendment double-jeopardy objection was overruled, and he was convicted of first-degree murder and sentenced to death (\href{https://supreme.justia.com/cases/federal/us/302/319/}{Justia}). Palko appealed to the Connecticut Supreme Court of Errors, which affirmed the conviction and death sentence in \textit{State v. Palko}, 122 Conn. 529, 191 A. 320 (1937), adhering to its 1894 precedent upholding the state-appeal statute in \textit{State v. Lee}, 65 Conn. 265, 30 A. 1110 (1894) (\href{https://supreme.justia.com/cases/federal/us/302/319/}{Justia}).

Palko appealed as of right to the U.S. Supreme Court under 28 U.S.C. \S~344 (\href{https://supreme.justia.com/cases/federal/us/302/319/}{Justia}). The Court heard oral argument on November 12, 1937, and issued its decision on December 6, 1937 (\href{https://www.oyez.org/cases/1900-1940/302us319}{Oyez}). After the affirmance, Frank Palka was executed in Connecticut's electric chair on April 12, 1938 (\href{https://www.oyez.org/cases/1900-1940/302us319}{Oyez}).

\section{Judicial Votes}

The Court voted \textbf{8--1} to affirm the Connecticut Supreme Court of Errors (\href{https://www.oyez.org/cases/1900-1940/302us319}{Oyez}; \href{https://ballotpedia.org/Palko_v._Connecticut}{Ballotpedia}):
\begin{itemize}
    \item \textbf{Majority (8):} Justice Benjamin N. Cardozo delivered the opinion of the Court, joined by Chief Justice Charles Evans Hughes and Justices James Clark McReynolds, Louis D. Brandeis, George Sutherland, Harlan Fiske Stone, Owen J. Roberts, and Hugo L. Black (\href{https://www.oyez.org/cases/1900-1940/302us319}{Oyez}; \href{https://ballotpedia.org/Palko_v._Connecticut}{Ballotpedia}). (Notably, Justice Black---newly appointed in 1937---joined Cardozo's opinion, though he would later champion total incorporation in \textit{Adamson v. California}, 332 U.S. 46 (1947).)
    \item \textbf{Dissent (1):} Justice Pierce Butler dissented without writing an opinion (\href{https://supreme.justia.com/cases/federal/us/302/319/}{Justia}; \href{https://ballotpedia.org/Palko_v._Connecticut}{Ballotpedia}).
\end{itemize}

\section{Holding}

\textbf{Judgment:} Affirmed (\href{https://supreme.justia.com/cases/federal/us/302/319/}{Justia}).

\textbf{Rule Announced:} The Fourteenth Amendment does not automatically incorporate all guarantees of the first eight amendments of the Bill of Rights against the states, under either the Due Process Clause or the Privileges or Immunities Clause (\href{https://supreme.justia.com/cases/federal/us/302/319/}{Justia}). Instead, a Bill of Rights guarantee is absorbed into the Fourteenth Amendment's Due Process Clause only if it is ``of the very essence of a scheme of ordered liberty''---meaning a ``principle of justice so rooted in the traditions and conscience of our people as to be ranked as fundamental'' such that ``neither liberty nor justice would exist if [it] were sacrificed'' (\href{https://supreme.justia.com/cases/federal/us/302/319/}{Justia}). Because a state-initiated retrial solely to correct substantial legal error in the same case does not subject the accused to ``a hardship so acute and shocking that our polity will not endure it,'' Connecticut's statute did not violate the Fourteenth Amendment (\href{https://supreme.justia.com/cases/federal/us/302/319/}{Justia}).

\textbf{Holding vs. Dicta:} The narrow holding sustained Conn. Gen. Stat. \S~6494 and held that the Fifth Amendment's Double Jeopardy Clause did not bar state retrial after an error-infected first trial---a holding the Supreme Court expressly overruled in \textit{Benton v. Maryland}, 395 U.S. 784 (1969) (\href{https://www.oyez.org/cases/1900-1940/302us319}{Oyez}). Cardozo explicitly reserved in dicta what the Constitution would require if a state tried an accused again ``after a trial free from error'' or brought a separate prosecution to wear the defendant down (\href{https://supreme.justia.com/cases/federal/us/302/319/}{Justia}). Equally consequential were Cardozo's broader doctrinal observations classifying Bill of Rights provisions: while his remarks that criminal jury trials and the privilege against self-incrimination were not essential to ordered liberty were later repudiated (\textit{Duncan v. Louisiana}, 391 U.S. 145 (1968); \textit{Malloy v. Hogan}, 378 U.S. 1 (1964)), his ``ordered liberty'' framework and his description of freedom of thought and speech as ``the matrix, the indispensable condition, of nearly every other form of freedom'' became foundational pillars of modern constitutional law (\href{https://www.law.cornell.edu/constitution-conan/amendment-14/section-1/due-process-clause-and-incorporation-early-doctrine}{Cornell LII Constitution Annotated}; \href{https://supreme.justia.com/cases/federal/us/302/319/}{Justia}).

\section{Analysis of Opinions}

\subsection{Majority Opinion (Justice Cardozo)}

Justice Cardozo began by assuming, without deciding, that under \textit{Kepner v. United States}, 195 U.S. 100 (1904), a government-initiated appeal and retrial in the same case would violate the Fifth Amendment's Double Jeopardy Clause in a federal prosecution (\href{https://supreme.justia.com/cases/federal/us/302/319/}{Justia}). Nevertheless, he warned against ``the tyranny of labels,'' rejecting Palko's broad thesis that ``whatever is forbidden by the Fifth Amendment is forbidden by the Fourteenth also'' and that the Fourteenth Amendment mechanically incorporates Amendments I through VIII (\href{https://supreme.justia.com/cases/federal/us/302/319/}{Justia}).

To demonstrate that ``there is no such general rule,'' Cardozo surveyed two lines of precedent (\href{https://supreme.justia.com/cases/federal/us/302/319/}{Justia}). On the non-incorporated side, the Court had previously held that states could dispense with grand jury indictments (\textit{Hurtado v. California}, 110 U.S. 516 (1884)), modify or abolish criminal and civil jury trials (\textit{Maxwell v. Dow}, 176 U.S. 581 (1900); \textit{Walker v. Sauvinet}, 92 U.S. 90 (1876)), and eliminate the privilege against self-incrimination (\textit{Twining v. New Jersey}, 211 U.S. 78 (1908)) (\href{https://supreme.justia.com/cases/federal/us/302/319/}{Justia}). On the incorporated side, the Court had held that the Fourteenth Amendment absorbed First Amendment freedoms of speech, press, free exercise of religion, and peaceable assembly (\textit{Near v. Minnesota}, 283 U.S. 697 (1931); \textit{De Jonge v. Oregon}, 299 U.S. 353 (1937)), as well as the right to counsel when necessary for a real hearing in a capital case (\textit{Powell v. Alabama}, 287 U.S. 45 (1932)) (\href{https://supreme.justia.com/cases/federal/us/302/319/}{Justia}).

Cardozo synthesized these cases under a single ``rationalizing principle'': rights are absorbed into the Fourteenth Amendment not because they appear in the Bill of Rights, but because they are ``implicit in the concept of ordered liberty'' such that ``neither liberty nor justice would exist if they were sacrificed'' (\href{https://supreme.justia.com/cases/federal/us/302/319/}{Justia}). Procedural mechanisms like grand juries, petit juries, and the privilege against self-incrimination---though valuable---were not indispensable to a ``fair and enlightened system of justice,'' as evidenced by Continental European systems (so long as states barred torture under \textit{Brown v. Mississippi}, 297 U.S. 278 (1936)) (\href{https://supreme.justia.com/cases/federal/us/302/319/}{Justia}). By contrast, freedom of thought and speech occupied a ``different plane of social and moral values'' as the ``matrix'' of nearly every other freedom (\href{https://supreme.justia.com/cases/federal/us/302/319/}{Justia}).

Applying this test to Conn. Gen. Stat. \S~6494, Cardozo reasoned that Connecticut was not attempting to wear Palko down through repeated prosecutions after an error-free trial; rather, the State asked only for ``a trial free from the corrosion of substantial legal error,'' exercising a reciprocal appellate privilege supervised by the trial judge (\href{https://supreme.justia.com/cases/federal/us/302/319/}{Justia}). ``The edifice of justice stands,'' Cardozo concluded, ``its symmetry, to many, greater than before'' (\href{https://supreme.justia.com/cases/federal/us/302/319/}{Justia}). Finally, Cardozo summarily rejected Palko's Fourteenth Amendment Privileges or Immunities Clause claim in two sentences, citing \textit{Maxwell v. Dow}, 176 U.S. at 584 (\href{https://supreme.justia.com/cases/federal/us/302/319/}{Justia}).

\subsection{Dissent (Justice Butler)}

Justice Pierce Butler entered a single-line dissent (``Mr. Justice Butler dissents'') without publishing an opinion explaining his reasoning (\href{https://supreme.justia.com/cases/federal/us/302/319/}{Justia}). Consequently, the record leaves no formal doctrinal statement from the lone dissenter, though his vote signaled disagreement with allowing Connecticut to retry Palko for his life after a jury had already convicted him of the lesser offense of second-degree murder.

\section{Examples: Future Applications}

\begin{enumerate}
    \item \textbf{Same-Side (Not Incorporated under \textit{Palko}'s Framework): State Prosecution by Prosecutor's Information Rather Than Grand Jury Indictment.} A state charges a defendant with felony burglary via a prosecutor's information following a preliminary hearing rather than a Fifth Amendment grand jury indictment. Under \textit{Palko}'s ``ordered liberty'' standard---and still today---the Fifth Amendment's Grand Jury Clause is not incorporated against the states because a fair, accurate criminal adjudication remains possible without a grand jury (\href{https://supreme.justia.com/cases/federal/us/302/319/}{Justia}; \href{https://www.cardozolawreview.com/interrogating-the-nonincorporation-of-the-grand-jury-clause/}{Cardozo Law Review}).

    \item \textbf{Same-Side (Not Incorporated under \textit{Palko}'s Framework): State Abolition of Civil Jury Trials in Tort Suits Exceeding \$20.} A state legislature replaces jury trials with bench trials for all civil contract and tort disputes. Under \textit{Palko}, which explicitly cited \textit{Walker v. Sauvinet}, 92 U.S. 90 (1876), the Seventh Amendment's civil jury guarantee is not essential to a scheme of ordered liberty and remains unincorporated against the states (\href{https://supreme.justia.com/cases/federal/us/302/319/}{Justia}).

    \item \textbf{Opposite-Side (Incorporated under \textit{Palko}'s Framework): State Prior Restraint on Political Criticism.} A state enacts a statute empowering local officials to enjoin newspapers from publishing articles critical of state judges. Under \textit{Palko}, freedom of thought, speech, and the press sit at the very apex of ``ordered liberty'' as the ``matrix'' of nearly every other freedom, so the First Amendment fully binds the state through the Fourteenth Amendment (\href{https://supreme.justia.com/cases/federal/us/302/319/}{Justia}).

    \item \textbf{Opposite-Side (Violates Due Process Even Under \textit{Palko}): State Retrial After an Error-Free Acquittal.} A state legislature authorizes prosecutors to retry a defendant a second time solely because the state disagrees with a jury's acquittal following a trial free from legal error. Even within \textit{Palko}'s own reasoning, Cardozo expressly distinguished error-free retrials aimed at wearing the accused down through ``a multitude of cases with accumulated trials'' as a potential ``hardship so acute and shocking'' as to violate fundamental fairness---and after \textit{Benton v. Maryland}, 395 U.S. 784 (1969), such retrials are categorically barred by the incorporated Double Jeopardy Clause (\href{https://supreme.justia.com/cases/federal/us/302/319/}{Justia}; \href{https://www.oyez.org/cases/1900-1940/302us319}{Oyez}).

    \item \textbf{Fence-Sitter (Shifted Across Eras): State Criminal Procedure Protections Such as Non-Unanimous Juries or Excessive Fines.} Under \textit{Palko}'s strict, universalist inquiry---asking whether \textit{any} ``fair and enlightened system of justice'' could exist without a rule---states enjoyed wide procedural leeway, and Cardozo expressly placed criminal jury trials and self-incrimination outside ordered liberty (\href{https://supreme.justia.com/cases/federal/us/302/319/}{Justia}). Under the modern selective incorporation standard refined in \textit{Duncan v. Louisiana}, 391 U.S. 145 (1968), and \textit{McDonald v. City of Chicago}, 561 U.S. 742 (2010)---which asks instead whether a right is fundamental to \textit{our} American scheme of ordered liberty and deeply rooted in this Nation's history and tradition---the Court incorporated jury trials, double jeopardy, excessive fines (\textit{Timbs v. Indiana}, 586 U.S. 146 (2019)), and unanimous felony juries (\textit{Ramos v. Louisiana}, 590 U.S. 83 (2020)), making borderline procedural safeguards depend heavily on whether a court applies \textit{Palko}'s abstract universal test or \textit{Duncan}/\textit{McDonald}'s American-tradition test (\href{https://supreme.justia.com/cases/federal/us/561/742/}{Justia McDonald}; \href{https://www.law.cornell.edu/constitution-conan/amendment-14/modern-doctrine-on-selective-incorporation-of-bill-of-rights}{Cornell LII Constitution Annotated}).
\end{enumerate}

\section{Critique}

\textbf{Originalist and Textualist Perspectives:} Originalist and textualist critics attack \textit{Palko} from two opposite directions. First, ``total incorporation'' and Privileges or Immunities originalists---beginning with Justice Hugo Black's landmark dissent in \textit{Adamson v. California}, 332 U.S. 46 (1947), and later developed by scholars such as Akhil Reed Amar and Justice Thomas's concurrence in \textit{McDonald v. City of Chicago}, 561 U.S. 742 (2010)---argue that \textit{Palko} erred both by relying on the Due Process Clause rather than the Privileges or Immunities Clause and by empowering judges to pick and choose which enumerated rights count as ``ordered liberty'' (\href{https://www.cardozolawreview.com/interrogating-the-nonincorporation-of-the-grand-jury-clause/}{Cardozo Law Review}; \href{https://supreme.justia.com/cases/federal/us/561/742/}{Justia McDonald}). Second, anti-incorporation or strict-federalism critics argue that \textit{Palko} lacked textual fidelity to the Fourteenth Amendment and opened the door to subjective judicial policymaking untethered from constitutional text (\href{https://lonang.com/commentaries/conlaw/bill-of-rights/the-incorporation-doctrine/}{Lonang Institute}; \href{https://constitutingamerica.org/palko-v-connecticut-1937-guest-essayist-robert-lowry-clinton/}{Constituting America}).

\textbf{Progressive and Doctrinal Perspectives:} Doctrinal and civil-liberties scholars criticize \textit{Palko}'s abstract ``civilized system of justice'' benchmark as unduly stingy toward criminal defendants. As Louis Henkin and subsequent Warren Court opinions observed, asking whether Continental Europe could dispense with a right ignored the specific adversarial structure of Anglo-American criminal justice, where double jeopardy, jury trials, and the privilege against self-incrimination are vital bulwarks against state overreach (\href{https://openyls.law.yale.edu/server/api/core/bitstreams/8b7df10e-f7aa-4233-ba2b-2ca63a92071a/content}{Yale Law Journal}). Consequently, during the 1960s, the Court retained \textit{Palko}'s selective-incorporation structure while replacing Cardozo's abstract universalism with a focus on whether a guarantee is fundamental to the \textit{American} scheme of justice---culminating in \textit{Benton v. Maryland}, 395 U.S. 784 (1969), which rejected \textit{Palko}'s ``watered-down'' approach to double jeopardy and expressly overruled it (\href{https://www.law.cornell.edu/constitution-conan/amendment-14/modern-doctrine-on-selective-incorporation-of-bill-of-rights}{Cornell LII Constitution Annotated}; \href{https://www.oyez.org/cases/1900-1940/302us319}{Oyez}).

\textbf{Underlying Competing Values:} \textit{Palko} sits at the intersection of federalism and nationalized individual rights. Cardozo sought to preserve states as flexible laboratories of criminal procedure without sacrificing core democratic freedoms like speech and fair hearings. Yet the decision's tension---and its tragic human outcome, ending in Palko's execution after a state-initiated retrial---illustrates the enduring conflict between judicial deference to state experimentation and uniform national protection of the Bill of Rights.

\section{Key Quotations}

\begin{quote}
``His thesis is even broader. Whatever would be a violation of the original bill of rights (Amendments I to VIII) if done by the federal government is now equally unlawful by force of the Fourteenth Amendment if done by a state. There is no such general rule.''
\end{quote}
\noindent --- Justice Cardozo, Opinion of the Court, 302 U.S. at 323 (\href{https://supreme.justia.com/cases/federal/us/302/319/}{Justia}; \href{https://www.courtlistener.com/opinion/102879/palko-v-connecticut/}{CourtListener}).

\begin{quote}
``In these and other situations, immunities that are valid as against the federal government by force of the specific pledges of particular amendments have been found to be implicit in the concept of ordered liberty, and thus, through the Fourteenth Amendment, become valid as against the states.''
\end{quote}
\noindent --- Justice Cardozo, Opinion of the Court, 302 U.S. at 324--325 (\href{https://supreme.justia.com/cases/federal/us/302/319/}{Justia}).

\begin{quote}
``The right to trial by jury and the immunity from prosecution except as the result of an indictment may have value and importance. Even so, they are not of the very essence of a scheme of ordered liberty. To abolish them is not to violate a `principle of justice so rooted in the traditions and conscience of our people as to be ranked as fundamental.' ''
\end{quote}
\noindent --- Justice Cardozo, Opinion of the Court, 302 U.S. at 325 (\href{https://supreme.justia.com/cases/federal/us/302/319/}{Justia}).

\begin{quote}
``If the Fourteenth Amendment has absorbed them, the process of absorption has had its source in the belief that neither liberty nor justice would exist if they were sacrificed. This is true, for illustration, of freedom of thought, and speech. Of that freedom one may say that it is the matrix, the indispensable condition, of nearly every other form of freedom.''
\end{quote}
\noindent --- Justice Cardozo, Opinion of the Court, 302 U.S. at 326--327 (\href{https://supreme.justia.com/cases/federal/us/302/319/}{Justia}; \href{https://openjurist.org/302/us/319/palko-v-state-of-connecticut}{OpenJurist}).

\begin{quote}
``Is that kind of double jeopardy to which the statute has subjected him a hardship so acute and shocking that our polity will not endure it? Does it violate those `fundamental principles of liberty and justice which lie at the base of all our civil and political institutions'? The answer surely must be `no.' ''
\end{quote}
\noindent --- Justice Cardozo, Opinion of the Court, 302 U.S. at 328 (\href{https://supreme.justia.com/cases/federal/us/302/319/}{Justia}).

\end{document}
